% answers_short.tex —— 2011 年·天津卷·理科 · 简短答案速查
% 由 tools/make_short_answers.py 自 answers.tex 的 \gksolhead/\gkanswer 提取，
% 并与 2011/2011_数学真题_天津理_答案要点.md 逐题比对后生成，勿手改（改源 MD 后重跑）。

\begin{gknotice}
\textbf{使用说明}：本节只给最终答案便于快速核对，完整推导见解析版。答案为编写者独立推导，\textbf{非官方标准答案}。
\end{gknotice}

\gkpart{选择题}
\begin{center}\begin{tabular}{ccccccccc}
\toprule
题号 & 1 & 2 & 3 & 4 & 5 & 6 & 7 & 8 \\
\midrule
答案 & B & A & B & D & C & D & C & B \\
\bottomrule
\end{tabular}\end{center}

\gkpart{填空题}
\begin{itemize}[leftmargin=2.6em]
  \item[\textbf{9.}] $12$。
  \item[\textbf{10.}] $6+\pi$。
  \item[\textbf{11.}] $\sqrt2$。
  \item[\textbf{12.}] $\dfrac{\sqrt7}{2}$。
  \item[\textbf{13.}] $[-2,5]$。
  \item[\textbf{14.}] $5$。
\end{itemize}

\gkpart{解答题}
\begin{itemize}[leftmargin=2.6em]
  \item[\textbf{15.}] （1）定义域 $\left\{x\ \middle|\ x\neq\dfrac{\pi}{8}+\dfrac{k\pi}{2},\ k\in\mathbf{Z}\right\}$，最小正周期 $\dfrac{\pi}{2}$；（2）$\alpha=\dfrac{\pi}{12}$。
  \item[\textbf{16.}] （1）① $\dfrac15$；② $\dfrac{7}{10}$；（2）分布列见解析，$E(X)=\dfrac75$。
  \item[\textbf{17.}] （1）$\dfrac{\sqrt2}{3}$；（2）$\dfrac{3\sqrt5}{7}$；（3）$\dfrac{\sqrt{10}}{4}$。
  \item[\textbf{18.}] （1）$e=\dfrac12$；（2）$6x^2-\dfrac{16\sqrt3}{3}xy=5$（$x-\dfrac{y}{\sqrt3}>0$ 的部分）。
  \item[\textbf{19.}] （1）在 $\left(0,\dfrac{1}{\sqrt{2a}}\right)$ 上单调递增，在 $\left(\dfrac{1}{\sqrt{2a}},+\infty\right)$ 上单调递减；（2）（3）证明见解析。
  \item[\textbf{20.}] （1）$a_3=-3$，$a_4=-5$，$a_5=4$；（2）证明见解析，$c_n=(-1)^n$；（3）证明见解析。
\end{itemize}
